% Structure follows the tutor feedback: discovery -> soundness -> metric goal ->
% refinement -> conformance checking -> manual edits -> final model/BPMN ->
% branch probabilities -> decision rules. Tables/figures are factual; every
% '% -' line is a keyword note for the author's own text (no generated prose).

% ---------------------------------------------------------------- Discovery
% TODO (author): compare Alpha, Heuristics, Inductive, Genetic (docs/DISCOVERY.md,
% docs/ALPHA_MINER.md, docs/HEURISTICS_SWEEP.md, docs/GENETIC_BUDGET.md).
% - same input for all: completion projection, 31,509 cases, 475,306 events, 24 activities
% - Alpha \cite{vdaalst2004workflow}: WF-net by construction (1 source, 1 sink), but unbounded,
%   >=10 certified dead transitions, certified 0 cases replayable (short loops, noise, no silent transitions)
% - Heuristics: native output = heuristics net (App. figure), 16 end activities -> 70 silent
%   transitions in the Petri net; unbounded at dependency 0.5/0.7/0.9/0.99; >=12 dead; <=1.30% replayable
% - Genetic \cite{vdaalst2005genetic}: native objective = 0.4 trace fitness + 0.6 fitting share only;
%   2 of 6 generations (early stop); unbounded; final marking certified unreachable
% - Inductive \cite{leemans2013constructive}: sound by construction (block structure), 100% replay,
%   but Align-ET precision 0.293 -> starting point for refinement
% - HM token values not reproducible (+-0.03 between runs, hash-order silent-transition choice)
% Full table, all four criteria: App. Table~\ref{tab:model-comparison}; soundness: Table~\ref{tab:soundness}.
% Soundness notions \cite{vdaalst1997verification,vdaalst2011soundness}; Woflan cross-check \cite{verbeek2001woflan}.
% Factual output; source results/tables/quality_metrics.json
% Source SHA256: 316fda16a0f0877b7f1fb342541cb1ced289ddd6c0280a9320852b2a572bff1b
% Additional source results/tables/dead_transition_audit.json; SHA256 b5fdf629b415cf7ac3b4ac4f2e459195e0f68ce4517471c8f26876e10509be62
\begin{table}[htbp]
\centering
\small
\begin{tabular}{lrr}
\toprule
Metric & IM & HM \\
\midrule
Token replay fitness & 1.0000 & 0.8334 \\
Token mean trace fitness & 1.0000 & 0.8240 \\
Token fitting cases (\%) & 100.00 & 0.00 \\
Alignment fitness & 0.9998 & -- \\
Alignment mean trace fitness & 1.0000 & -- \\
Alignment fitting cases (\%) & 100.00 & -- \\
Replay upper bound (\%) & 100.00 & 1.30 \\
Token precision & 0.2926 & 0.9978 \\
Alignment precision & 0.2926 & 0.9975 \\
Unfit Align-ET prefixes (\%) & 0.00 & 84.35 \\
Token generalization & 0.9757 & 0.8611 \\
Arc-degree simplicity & 0.6170 & 0.4815 \\
\midrule
Nodes (own) & 87 & 143 \\
Arcs (own) & 114 & 220 \\
Extended Cardoso (own) & 47 & 108 \\
Extended Cyclomatic (own) & 1552 & -- \\
PM4Py Cyclomatic$^*$ & 2397 & -- \\
\bottomrule
\end{tabular}
\par\smallskip
\begin{minipage}{\linewidth}\scriptsize
$--$: unavailable; HM Cyclomatic: unbounded state space.\\
$^*$PM4Py 2.7.23.11: edge-name targets; own ECyM: marking targets.\\
Complexity: smaller values; quality scores and coverage: larger values.\\
Unfit-prefix share: smaller values; weighted nonempty prefixes.\\
Replay upper bound: unreachable-activity certificates.\\
Metric context~\cite{carmona2018conformance}; size~\cite{pedro2015simplification}; ECaM/ECyM~\cite{lassen2009complexity}.
\end{minipage}
\caption{Baseline metrics: 31,509 cases, 475,306 completed events.}
\label{tab:baseline-quality}
\end{table}
% Author: explain the metric definitions, model limitations, and the PM4Py Cyclomatic discrepancy.
% Definitions 13/15: Lassen and van der Aalst (2009), doi:10.1016/j.infsof.2008.08.005.
% Size: de San Pedro et al. (2015), doi:10.1007/978-3-319-23063-4_30.


% ---------------------------------------------------------------- Metric goal
% TODO (author): which metric to optimise and why (docs/MINER_PARAMETERS.md, D13).
% - goal = simulation base model
% - fitness as constraint: sound + ~80% exact whole-case replay (band 78--82%)
% - precision = optimisation target: unobserved behaviour -> unrealistic simulated cases
% - generalization, simplicity = guard rails (no notable loss)
% - counter-example: HM Align-ET ~0.998 although <=1.30% of cases can fit -> precision alone misleading
% Literature: \cite{carmona2018conformance,lassen2009complexity,pedro2015simplification}.

% ---------------------------------------------------------------- Refinement
% TODO (author): algorithm choice IMf (docs/REFINEMENT.md, docs/DESIGN_DECISIONS.md D05, D12--D16).
% - IMf \cite{leemans2014infrequent}: keeps soundness + block structure, noise threshold = precision lever
% - 40 candidates: 19 full-log noise values, 21 frequent-variant candidates; all evaluated on all cases
% - chosen: frequent 60% variants, noise 0.03 -> 80.45% exact replay, Align-ET 0.598
% - external comparisons: \cite{rodrigues2017stairway,povalyaeva2017density,blevi2017kpmg}
% Factual output; final_metrics.json SHA256 6e2dbf09f0630f4d55ffc61b9f2041114d15cd54740dc0000adfef2b0271ac6f
\begin{table}[htbp]
\centering
\small
\caption{Refined model versus IM baseline: same 31,509 cases; exact case replay from the complete marking graph.}
\label{tab:refined-quality}
\begin{tabular}{@{}lrr@{}}
\toprule
Metric & IM & Refined \\
\midrule
Exact complete-case replay (\%) & 100 & 80.4532 \\
Token replay fitness & 1.0000 & 0.9875 \\
Token fitting cases (\%) & 100.0000 & 81.6243 \\
Token precision & 0.2926 & 0.5952 \\
Align-ET precision & 0.2926 & 0.5977 \\
Token generalization & 0.9757 & 0.9912 \\
Arc-degree simplicity & 0.6170 & 0.7008 \\
Nodes & 87 & 89 \\
Arcs & 114 & 108 \\
Extended Cardoso & 47 & 52 \\
Extended Cyclomatic & 1,552 & 551 \\
\bottomrule
\end{tabular}
\end{table}


% ---------------------------------------------------------------- Conformance checking after mining
% TODO (author): conformance checking after mining as basis for manual edits (docs/MANUAL_EDITS.md).
% - precision leaks = Align-ET escaping edges (weighted prefixes); fitness leaks = alignment log moves
% - largest leaks: A_Complete parallel to whole offer loop; O_Sent self-loop; validation AND-block
% - 'Edit' column = edit targeting this leak; '--' = not addressed
\begin{table*}[t]
\centering
\small
\caption{Conformance diagnostics of the refined IMf model (before manual edits) and the edit addressing each leak; all 31,509 cases.}
\label{tab:conformance-edits}
\begin{tabular}{llp{0.40\textwidth}rl}
\hline
Leak & Diagnostic & Where & Weight & Edit \\
\hline
precision & escaping edge & O\_Create Offer $\rightarrow$ A\_Complete & 35,006 & R1 \\
 &  & O\_Created $\rightarrow$ A\_Complete & 35,006 & R1 \\
 &  & O\_Sent (mail and online) $\rightarrow$ O\_Sent (mail and online) & 33,899 & R2 \\
 &  & O\_Sent (mail and online) $\rightarrow$ O\_Create Offer & 33,482 & -- \\
 &  & A\_Validating $\rightarrow$ A\_Validating & 33,294 & R3 \\
 &  & A\_Validating $\rightarrow$ W\_Call incomplete files & 33,294 & R3 \\
 &  & A\_Accepted $\rightarrow$ A\_Complete & 31,483 & R1 \\
 &  & A\_Complete $\rightarrow$ O\_Returned & 30,349 & R3 \\
 &  & A\_Complete $\rightarrow$ W\_Validate application & 30,349 & -- \\
 &  & A\_Complete $\rightarrow$ A\_Incomplete & 30,349 & R3 \\
 &  & A\_Validating $\rightarrow$ O\_Accepted & 20,966 & -- \\
 &  & A\_Incomplete $\rightarrow$ W\_Validate application & 19,626 & -- \\
\hline
fitness & log move & O\_Sent (online only) & 2,026 & G1 \\
 &  & O\_Cancelled & 1,964 & G2, R4 \\
 &  & O\_Returned & 1,904 & -- \\
 &  & O\_Created & 1,575 & -- \\
 &  & O\_Create Offer & 1,575 & -- \\
 &  & O\_Refused & 991 & G3 \\
\hline
\end{tabular}

\end{table*}

% ---------------------------------------------------------------- Manual edits
% TODO (author): which flow was changed manually and why (docs/MANUAL_EDITS.md).
% - edits on the process tree, not single arcs: soundness + block structure preserved; Petri-net delta logged
% - R = restricting (precision up, replay down), G = generalizing (replay up, precision slightly down)
% - all 128 combinations screened; 40 in band; selection rule fixed in advance
% - R1: strongest evidence (0 counter-cases), +0.058 precision, but duplicates offer block ->
%   alignment generalization 0.972 vs 0.991 -> fails guard (all 14 R1 combinations)
% - R2: 8.83% of cases repeat O_Sent -> leaves band
% - selected G3+R3+R4: Align-ET 0.598 -> 0.697, exact replay 80.45% -> 79.39%, generalization -0.001
% - guard stage, all checked candidates: App. Table~\ref{tab:guard-stage}
\begin{table*}[t]
\centering
\small
\caption{Single manual edits of the refined IMf process tree: log evidence and effect; all 31,509 cases.}
\label{tab:single-edits}
\begin{tabular}{llp{0.36\textwidth}rrrr}
\hline
Edit & Type & Log evidence (all 31,509 cases) & Exact replay \% & $\Delta$pp & Align-ET & Nodes \\
\hline
base & -- & refined IMf (frequent\_0.60\_imf\_0.03) & 80.45 & -- & 0.598 & 89 \\
R1 & restricting & A\_Complete after first O\_Sent: 31,362 (99.53\%); before: 0 & 80.42 & -0.03 & 0.656 & 111 \\
R2 & restricting & direct O\_Sent repetition in 2,781 (8.83\%) & 74.17 & -6.28 & 0.624 & 84 \\
R3 & restricting & O\_Returned / A\_Incomplete before first A\_Validating: 0 / 0 & 77.90 & -2.55 & 0.697 & 84 \\
R4 & restricting & O\_Cancelled after A\_Denied: 195; after A\_Cancelled: 10,431 & 80.45 & +0.00 & 0.598 & 92 \\
G1 & generalizing & O\_Sent (online only) in 1,800 (5.71\%) & 81.88 & +1.43 & 0.585 & 90 \\
G2 & generalizing & A\_Cancelled after A\_Validating: 1,014 (3.22\%) & 81.52 & +1.07 & 0.595 & 90 \\
G3 & generalizing & $\geq$2 O\_Refused: 777 (2.47\%) & 82.06 & +1.61 & 0.597 & 92 \\
\hline
\end{tabular}

\end{table*}

% ---------------------------------------------------------------- Final model
% TODO (author): final model quality, trade-offs, generalization evidence.
% - four criteria IM -> IMf -> final: precision up, replay constraint kept, generalization stable
% - held-out: IMf rediscovered on 70% train cases only = identical tree; test replay 80.37% vs train 80.49%
% - caveat: edits screened on all cases
\begin{table}[t]
\centering
\small
\caption{Four quality criteria along the refinement path; all 31,509 cases.}
\label{tab:final-quality}
\resizebox{\columnwidth}{!}{\begin{tabular}{llrrr}
\hline
Criterion & Measure & IM & IMf & Final \\
\hline
Fitness & exact replay \% & 100.00 & 80.45 & 80.43 \\
 & alignment & 1.000 & 0.977 & 0.979 \\
Precision & Align-ET & 0.293 & 0.598 & 0.695 \\
Generalization & alignment & 0.976 & 0.991 & 0.988 \\
Simplicity & arc degree & 0.617 & 0.701 & 0.710 \\
 & nodes & 87 & 89 & 98 \\
 & Ext.\ Cardoso & 47 & 52 & 55 \\
 & ECyM & 1,552 & 551 & 166 \\
\hline
\end{tabular}
}
\end{table}
\begin{table}[t]
\centering
\small
\caption{Exact whole-case replay on the case-wise 70/30 split (seed 20261018).}
\label{tab:holdout}
\resizebox{\columnwidth}{!}{\begin{tabular}{lrrrr}
\hline
Model & Train $n$ & Train exact \% & Test $n$ & Test exact \% \\
\hline
IMf (before edits) & 22,056 & 80.49 & 9,453 & 80.37 \\
Final (after edits) & 22,056 & 79.45 & 9,453 & 79.26 \\
\hline
Train-only IMf & \multicolumn{4}{l}{identical tree (72 variants, 13,258 cases)} \\
\hline
\end{tabular}
}
\end{table}

% ---------------------------------------------------------------- BPMN block structure
% TODO (author): BPMN from the process tree, block structure (docs/MANUAL_EDITS.md, Leemans 2013).
% - every operator = single-entry/single-exit block; X and loop -> XOR split/merge pair, + -> AND pair
% - block-structured process tree -> sound WF-net; BPMN language-equivalent to the Petri net
% - full BPMN: App. Figure~\ref{fig:final-edited-bpmn}
\begin{table*}[t]
\centering
\small
\caption{Block structure of the final model: process-tree operators and the BPMN constructs they become.}
\label{tab:block-structure}
\begin{tabular}{llr}
\hline
Process-tree block & BPMN construct & Count \\
\hline
sequence $\rightarrow$ & sequence flows, no gateway & 9 \\
exclusive choice $\times$ & XOR split + XOR merge pair & 9 \\
loop $\circlearrowleft$ & XOR merge (loop entry) + XOR split (exit or redo) pair & 8 \\
parallel $\wedge$ & AND split + AND join pair & 2 \\
activity leaf & task & 23 \\
tau leaf & no task; skip or redo path between gateways & 15 \\
\hline
\multicolumn{2}{l}{BPMN: tasks / XOR split+merge / AND split+join} & 23 / 17+17 / 2+2 \\
\hline
\end{tabular}

\end{table*}

% ---------------------------------------------------------------- Decision mining
% TODO (author): decision mining in two steps (docs/DECISION_MINING.md).
% - step 1: branch probabilities per XOR split from optimal alignments, no data (App. Figure~\ref{fig:branch-probabilities})
% - loop splits: probability per iteration, not per case
% - step 2: Algorithm 1 of \cite{deleoni2013dataaware}; CART instead of C4.5; case-wise 70/30 split
% - 17 XOR splits (9 choices + 8 loops), 17 merges: no choice, no classifier; 424,194 aligned decisions
% - 9 of 17 trees predict only the majority branch (balanced accuracy 0.5);
%   7 beat the majority baseline by >1 pp (D02, D03, D07, D10, D13, D14, D16)
% - structural vs guarded replay: 79.39% -> 10.33% (test 10.07%); cause: majority trees forbid minority branches
% - simulation: probabilities where a tree does not beat the majority baseline (open design choice)
\begin{table}[t]
\centering
\caption{Exact whole-case replay of the same BPMN structure, before and after enforcing learned guards.}
\label{tab:guarded-coverage}
\begin{tabular}{lrrr}
\hline
Cases & $n$ & Structure (\%) & Guards (\%) \\
\hline
Train & 22,056 & 80.48 & 9.80 \\
Test & 9,453 & 80.32 & 9.46 \\
All & 31,509 & 80.43 & 9.70 \\
\hline
\end{tabular}
\end{table}


% Design rationale map for your own discussion: docs/DESIGN_DECISIONS.md.
% D01--D05: input projection, order, weighting and common Petri-net semantics.
% D06--D11: IM/HM baseline defaults, Genetic budget/seeds and ordering check.
% D12--D18: exact coverage vs fitness; precision under constraints; complexity
% disagreements, frequent-variant training, rare-outcome bias and exclusions.
% D19--D22: alignment ambiguity, availability/leakage, grouped split, CART
% adaptation and the guard-coverage loss. D23--D25: execution/export/report choices.
% D26: unique PNML IDs, lossless source-graph correction and independent native-objective replay.
